\section{Evaluation}
\label{sec:evaluation}

We ask: (E1) does the system \emph{really} verify SEV-SNP attestation on
production confidential hardware? (E2) can real AI workloads run from governed
pods? (E3) does scheduling behave correctly across multiple confidential nodes
and invalid evidence states? (E4) does the deployed system withstand an
adversarial campaign? (E5) how does it compare, quantitatively, to a strong
scheduling-gate baseline --- our central result? (E6) is each enforcement layer,
and PreBind in particular, necessary? (E7) is the placement decision
cryptographically bound and offline-verifiable? (E8) what is the B1--B5 overhead?
and (E9) is the scheduler itself minimally privileged? Every number is backed by
a raw artifact under \texttt{article1/results/raw/}; each result states its
environment.

\subsection{Experimental environment}
Table~\ref{tab:env} records the exact real environment; all figures come from
these raw Azure outputs, not from memory. A local test cluster is used
\emph{only} for CI, debug, and regression; no such number is used as
principal security or performance evidence.

\begin{table}[t]
  \centering
  \caption{Experimental environment (from raw Azure outputs).}
  \label{tab:env}
  \footnotesize
  \begin{tabular}{@{}lp{4.9cm}@{}}
    \toprule
    Cloud provider & Microsoft Azure (AKS) \\
    Region & \texttt{westus2} \\
    Confidential node pool & \texttt{conf}, 4 active nodes during final AKS runs \\
    Confidential VM SKU & \texttt{Standard\_DC8as\_v6} (AMD SEV-SNP CVM) \\
    Quota family & Standard DCasv6 Family vCPUs (32) \\
    vCPU (confidential active/max) & 32 active / 32 max \\
    System pool & 1$\times$ \texttt{Standard\_D2s\_v3} (non-confidential) \\
    Runtime class & \texttt{runc} (node-level SEV-SNP; not pod sandbox) \\
    Attestation service & Microsoft Azure Attestation (MAA) \\
    MAA endpoint & \url{sharedwus2.wus2.attest.azure.net} \\
    Attestation type claim & \texttt{sevsnpvm} \\
    Experiment date & 2026-07-06 \\
    \bottomrule
  \end{tabular}
\end{table}

\noindent On \texttt{Standard\_DC8as\_v6} the pods run under \texttt{runc}: the
node is a confidential VM (SEV-SNP), and attestation is \emph{node-level}. AKS
does not enable pod sandboxing / Kata on this SKU (no nested virtualization), so
we do not claim pod-level \texttt{kata-vm-isolation}.

\subsection{E1: Real SEV-SNP attestation}
The node-attestation-agent collected genuine MAA tokens and the central verifier
appraised them. The final multi-node snapshot
(\texttt{evidence-initial.json}) contains four active
\texttt{Standard\_DC8as\_v6} confidential nodes, each with non-simulated
\texttt{evidenceMode=real}, \texttt{verificationStatus=Verified}, a
\texttt{maaTokenHash}, and a \texttt{claimsDigest}. The filtered summary
(\texttt{attestation-real-summary.json}) is used to exclude old/orphan evidence
objects from earlier autoscaler instances. The fact that the token passed RS256
verification against MAA's published JWKS establishes authentic hardware
attestation, not a replayed constant. Critically, the same verifier returns
\texttt{unverified}/\texttt{Failed} for an empty or malformed token (observed
during bring-up and enforced by unit tests
\texttt{TestUnavailableWhenKeysUnreachable}, \texttt{TestBadSignatureFails}):
the pipeline never fabricates a real result.

\subsection{E2: Governed AI workloads}
\label{sec:e2-ai}
Confidential AI inference is the motivating use case, so the AKS campaign runs
real, minimal AI workloads from attestation-governed pods rather than using AI as
a slogan. The workload CSV (\texttt{ai\_workloads.csv}) contains three passing
workloads: \texttt{openai-chat-minimal} (OpenAI Responses API, 3/3 requests,
median 2216.815\,ms), \texttt{openai-embedding-minimal} (OpenAI embeddings, 3/3,
median 705.118\,ms), and \texttt{local-cpu-minimal} (a deterministic local CPU
model, 5/5, median 0.001\,ms). All three pods are governed by a
\texttt{ConfidentialInferencePolicy}, receive \texttt{placement\_decision=allow},
and pass offline placement-token verification. We evaluate attested node
placement and binding to model/image digest identifiers; we do not claim model
confidentiality, OpenAI service confidentiality, GPU execution, or large-model
throughput.

\subsection{E3: Multi-node node qualification}
\label{sec:e3-multinode}
To avoid evaluating a scheduler on a single active confidential node, we scaled
the AKS confidential pool to four active \texttt{Standard\_DC8as\_v6} nodes under
the 32-vCPU DCasv6 quota. The snapshot used by
\texttt{run\_multinode\_node\_selection.sh} records four active confidential
nodes with real verified evidence. The experiment then exercises four cases
(\texttt{multinode\_node\_selection.csv}): a valid confidential node is selected
and bound with a verified placement token; a node whose real evidence is
temporarily marked revoked is not bound; a node whose real evidence is
temporarily aged beyond the freshness bound is not bound; and a non-confidential
system node selected by \texttt{nodeSelector} but lacking
\texttt{AttestationEvidence} is not bound. The revoked/expired modifications are
backed up and restored by the harness before exit.

\subsection{E4: Deployed adversarial campaign}
\label{sec:e2}
We executed, against the deployed AKS system, attacks A1--A10 plus the A5b
sub-case (eleven attack types), each from an adversarial principal, $N{=}30$
per attack. The merged raw matrix (\texttt{security\_attacks\_A1\_A11.csv})
contains 330 A1--A10/A5b attack instances (11 types $\times$ 30 runs) and all
are blocked; \texttt{security\_attack\_matrix.csv} reports Wilson 95\%
intervals per attack, including A5b. The rerun includes stale evidence (A2), revoked evidence
(A3), wrong TEE (A4), production-forbidden/simulated runtime use (A5/A9), missing
model digest (A5b), policy and revocation races between Filter and PreBind
(A6/A7), token tampering (A8), and forged raw reports / compromised node-agent
attempts (A10). A11 is reported separately as a GPU-scope fail-closed check: a
workload requiring confidential GPU attestation is denied when no such evidence
exists, but we do not claim GPU confidential-computing evaluation.

\subsection{E5: Strong baseline B4 vs.\ B5 (central result)}
\label{sec:e4}
This is the paper's central quantitative result. \textbf{B4} is a genuinely
strong, non-strawman baseline: an external controller places a
\texttt{schedulingGate} on the pod, verifies the \emph{same} real
\texttt{AttestationEvidence} (TEE, freshness, revocation), removes the gate, and
lets the default scheduler bind. It is exactly the design a careful engineer
would build with today's Kubernetes primitives. \textbf{B5} is our
attestation-aware scheduler with the fail-closed \texttt{PreBind} re-check and a
signed placement token. Over $N{=}30$ runs each (\texttt{b4\_vs\_b5.csv},
Fig.~\ref{fig:toctou}), B4 exhibits a gate-removal$\to$bind \textbf{TOCTOU window
(median $\approx$1.1\,s, up to $\sim$1.5\,s in this campaign)} --- the window in which a revocation
or policy change is not re-checked before bind --- and produces \emph{no}
offline-verifiable artifact. B5 eliminates this externally schedulable
gate-removal-to-bind window: the policy/evidence re-check runs inside
\texttt{PreBind}, and B5 emits a token that \texttt{verify-placement} validates
without cluster access. A residual in-process PreBind-to-Bind interval remains
theoretically possible, but it is bounded to a single local scheduling path and
does not hand control back to the API server between check and bind. The
difference is statistically decisive (Mann--Whitney $U{=}900$, $p<10^{-5}$;
Cliff's $\delta{=}1.0$, ``large''; \texttt{b4\_vs\_b5\_stats.csv}). The message is
threefold: schedulingGates are strong but leave an externally controllable
check-to-bind gap; PreBind removes that gap from the schedulable path; and the
placement token makes the outcome auditable --- properties B4 cannot provide.

\subsection{E6: Ablation --- is each layer necessary?}
\label{sec:e5}
We disable enforcement layers one at a time and re-run one representative attack
per added layer (\texttt{ablation.csv}, Fig.~\ref{fig:ablation}). With
\emph{policy-only} (L0, no scheduler enforcement) a forged-label pod is
\emph{admitted} (A1 bypasses): labels alone are not a control. Adding admission
(L1) blocks the runtime-class attack (A5). Adding the attestation scheduler (L2)
blocks the wrong-TEE-evidence attack (A4): the scheduler's Filter stage rejects
placement without valid, correctly-typed evidence. Adding the placement token
(L3) blocks token tampering (A8). Label Resistance specifically (forged-label
attack A1 blocked once the scheduler is present) is evidenced by the full-system
campaign rather than by this ablation, since A1 is blocked in 30/30 runs once
the scheduler is part of the deployed stack (\S\ref{sec:e2}, E4). The ablation
makes the central security claim falsifiable: removing a layer restores the
bypass for that layer's representative attack, so the block rate is
attributable to the mechanism, not to the choice of scenarios. The
intra-scheduler Filter-only vs.\ +PreBind ablation requires a build flag and is
noted as future work; E5 already isolates PreBind's effect quantitatively.

\subsection{E7: Identity binding \& verifiable placement}
\label{sec:e7}
The placement token binds \{pod UID, pod-spec hash, node identity, evidence hash,
policy hash\}. Using \texttt{verify-placement} offline
(\texttt{identity\_binding.csv}, Fig.~\ref{fig:idb}), the correct binding
verifies, and mutating \emph{any} bound field --- pod UID, pod-spec hash, node
identity, evidence hash, or policy hash --- makes verification fail (6/6 as
expected). A token minted for one pod cannot be replayed for another; a third
party can check ``where did my workload run?'' without cluster access.

\subsection{E8: B1--B5 overhead and scalability scope}
The AKS B1--B5 latency campaign contains two warmups and 30 measured runs per
baseline (\texttt{performance\_b1\_b5\_high\_resolution.csv}); all 150 measured
pods succeed, and the anti-quantization check reports a 0.0 ratio of exact
1000\,ms multiples. To avoid mixing instruments, Fig.~\ref{fig:perf-b1-b5}
plots only comparable client-observed scheduling-path latencies: for B1--B4 this
is the default scheduler's own reported total scheduling latency; because our
out-of-tree scheduler instead reports an in-process phase timer for this field
(see below), B5's comparable value is computed the same way an external
observer would measure any baseline --- pod-pending duration minus admission
latency --- rather than read from that field. The medians are
1170.5\,ms for B1 default scheduling, 1189.0\,ms for B2 label/nodeSelector,
1168.5\,ms for B3 RuntimeClass-only, 1128.0\,ms for B4 after the gate is
removed, and 1167.0\,ms for B5. Separately, B5's instrumented scheduler-internal
phase median is 216.232\,ms; this value comes from monotonic clock
instrumentation inside the attestation scheduler and is not plotted or compared
as a B1--B5 end-to-end latency. We therefore claim measured prototype overhead,
not a universal Kubernetes latency bound.
Scheduler-only scalability (KWOK) and local-test-cluster throughput are retained as
CI/debug/regression harnesses; we do not present them as principal cloud
security or performance results.

\begin{figure}[t]
  \centering
  \includegraphics[width=0.86\linewidth]{scheduling_latency_cdf.pdf}
  \caption{AKS B1--B5 scheduling-latency CDF using one comparable instrument:
  high-resolution client-observed scheduling-path timing. B5's separate
  scheduler-internal phase timing is reported only in text, not mixed into this
  CDF. Warmups are retained in raw data but excluded from this figure.}
  \label{fig:perf-b1-b5}
\end{figure}

\begin{figure}[t]
  \centering
  \includegraphics[width=0.86\linewidth]{toctou_window_b4_b5.pdf}
  \caption{Central result: the strong scheduling-gate baseline B4 leaves a
  $\sim$1.1\,s externally controllable gate-removal-to-bind window; our
  fail-closed PreBind design removes that window from the schedulable path
  ($N{=}30$; $p<10^{-5}$, Cliff's $\delta{=}1.0$).}
  \label{fig:toctou}
\end{figure}

\begin{figure}[t]
  \centering
  \includegraphics[width=0.86\linewidth]{ablation_security_latency.pdf}
  \caption{Ablation on real AKS: with policy only (L0) the forged-label attack
  bypasses; each added layer (admission, scheduler, token) blocks its attack ---
  the block rate is attributable to the mechanism.}
  \label{fig:ablation}
\end{figure}

\begin{figure}[t]
  \centering
  \includegraphics[width=0.86\linewidth]{identity_binding_matrix.pdf}
  \caption{Identity binding: the correct placement token verifies; mutating any
  bound field (pod UID, pod-spec/evidence/policy hash, node) is rejected offline.}
  \label{fig:idb}
\end{figure}

\begin{figure}[t]
  \centering
  \includegraphics[width=0.86\linewidth]{security_attack_heatmap.pdf}
  \caption{Deployed adversarial campaign on real SEV-SNP AKS. A1--A10 and A5b
  are each blocked in 30/30 runs; A11 is a single fail-closed GPU-scope control,
  not a GPU confidential-computing evaluation.}
  \label{fig:security}
\end{figure}

\subsection{E9: Scheduler RBAC minimality audit}
The scheduler is part of the TCB, so we audit its own privileges
(\S\ref{sec:scheduler-security}). Twelve \texttt{kubectl auth can-i} probes on the
real cluster (\texttt{scheduler\_security\_tests.csv}) confirm the nine forbidden
capabilities (forge/modify/delete evidence, create raw reports, patch node
labels, read arbitrary secrets, create webhook configs, wildcards) are denied and
the three required ones allowed. We deliberately scope this as an \emph{RBAC
minimality audit}; the complementary hardening (non-root, read-only root
filesystem, dropped capabilities, signing key in a Secret not readable via the
scheduler's own RBAC, fail-closed on a missing key) is described in
\S\ref{sec:scheduler-security}, and a fuller fake-scheduler / NetworkPolicy /
signed-image study is future work.

\subsection{Reproducibility}
Within our artifact, the campaign is a single command
(\texttt{run\_full\_campaign.sh}): preflight $\to$ provision $\to$ real-attestation
$\to$ E1--E9 $\to$ statistics/figures $\to$ verified teardown. All raw CSVs carry
\texttt{env}, \texttt{run\_id}, \texttt{seed}, \texttt{raw\_log\_path}; figures
come from \texttt{analyze\_all.py},
\texttt{update\_q1\_final\_artifacts.py}, and the AKS performance harness
(bootstrap 95\% CI, Mann--Whitney U). The artifact will be made publicly
available upon publication.
